\documentclass[10pt,a4paper]{amsart}
\usepackage{amsmath,amssymb,amsthm,graphicx}
% Retained minimal preamble for the Stacks Project chapter flat.tex (More on Flatness).
% Verbatim excerpt of the chapter preamble: Xy-pic setup (native preamble lines 14-18) and the native
% theorem declarations with the mathematical macros (native preamble lines 154-217). Copied mechanically;
% nothing added, removed or redefined. The wrapper class and the amsmath/amssymb/amsthm/graphicx base
% packages are supplied by the assembler, so the class/externaldocument part of the native preamble is omitted.
% For commutative diagrams we use Xy-pic
\usepackage[all]{xy}

% We use 2cell for 2-commutative diagrams.
\xyoption{2cell}
% Theorem environments.
%
\theoremstyle{plain}
\newtheorem{theorem}[subsection]{Theorem}
\newtheorem{proposition}[subsection]{Proposition}
\newtheorem{lemma}[subsection]{Lemma}

\theoremstyle{definition}
\newtheorem{definition}[subsection]{Definition}
\newtheorem{example}[subsection]{Example}
\newtheorem{exercise}[subsection]{Exercise}
\newtheorem{situation}[subsection]{Situation}

\theoremstyle{remark}
\newtheorem{remark}[subsection]{Remark}
\newtheorem{remarks}[subsection]{Remarks}

\numberwithin{equation}{subsection}

% Macros
%
\def\lim{\mathop{\mathrm{lim}}\nolimits}
\def\colim{\mathop{\mathrm{colim}}\nolimits}
\def\Spec{\mathop{\mathrm{Spec}}}
\def\Hom{\mathop{\mathrm{Hom}}\nolimits}
\def\Ext{\mathop{\mathrm{Ext}}\nolimits}
\def\SheafHom{\mathop{\mathcal{H}\!\mathit{om}}\nolimits}
\def\SheafExt{\mathop{\mathcal{E}\!\mathit{xt}}\nolimits}
\def\Sch{\mathit{Sch}}
\def\Mor{\mathop{\mathrm{Mor}}\nolimits}
\def\Ob{\mathop{\mathrm{Ob}}\nolimits}
\def\Sh{\mathop{\mathit{Sh}}\nolimits}
\def\NL{\mathop{N\!L}\nolimits}
\def\CH{\mathop{\mathrm{CH}}\nolimits}
\def\proetale{{pro\text{-}\acute{e}tale}}
\def\etale{{\acute{e}tale}}
\def\QCoh{\mathit{QCoh}}
\def\Ker{\mathop{\mathrm{Ker}}}
\def\Im{\mathop{\mathrm{Im}}}
\def\Coker{\mathop{\mathrm{Coker}}}
\def\Coim{\mathop{\mathrm{Coim}}}

% Boxtimes
%
\DeclareMathSymbol{\boxtimes}{\mathbin}{AMSa}{"02}

%
% Macros for moduli stacks/spaces
%
\def\QCohstack{\mathcal{QC}\!\mathit{oh}}
\def\Cohstack{\mathcal{C}\!\mathit{oh}}
\def\Spacesstack{\mathcal{S}\!\mathit{paces}}
\def\Quotfunctor{\mathrm{Quot}}
\def\Hilbfunctor{\mathrm{Hilb}}
\def\Curvesstack{\mathcal{C}\!\mathit{urves}}
\def\Polarizedstack{\mathcal{P}\!\mathit{olarized}}
\def\Complexesstack{\mathcal{C}\!\mathit{omplexes}}
% \Pic is the operator that assigns to X its picard group, usage \Pic(X)
% \Picardstack_{X/B} denotes the Picard stack of X over B
% \Picardfunctor_{X/B} denotes the Picard functor of X over B
\def\Pic{\mathop{\mathrm{Pic}}\nolimits}
\def\Picardstack{\mathcal{P}\!\mathit{ic}}
\def\Picardfunctor{\mathrm{Pic}}
\def\Deformationcategory{\mathcal{D}\!\mathit{ef}}

% --- appended 2026-10-12 (Chief): the project own definitions of the reference/slogan/history
% display environments, copied mechanically from the shared Stacks preamble lines 8-12
% (\usepackage{verbatim} plus the three \newenvironment lines, which the chapter preamble keeps
% verbatim from the project preamble). Nothing added, removed or redefined: these environments
% carry no mathematical content and the project itself typesets them as comments. Without these
% three lines a counted statement whose own text uses them fails with Environment slogan undefined.
\usepackage{verbatim}
\newenvironment{reference}{\comment}{\endcomment}
\newenvironment{slogan}{\comment}{\endcomment}
\newenvironment{history}{\comment}{\endcomment}
\usepackage{cleveref}
\crefname{originalreference}{original reference}{original references}
\Crefname{originalreference}{Original reference}{Original references}
\makeatletter
\newlabel{situation-existence-derived}{{0DIB}{source}{}{sourcepointer-1}{}}
\newlabel{situation-existence-derived@cref}{{[originalreference][0][]0DIB}{[1][0][]source}{}{}{}}
\makeatother
\crefname{originallemma}{original lemma}{original lemmas}
\Crefname{originallemma}{Original lemma}{Original lemmas}
\makeatletter
\newlabel{lemma-compute-what-it-should-be-derived}{{0DIC}{source}{}{sourcepointer-2}{}}
\newlabel{lemma-compute-what-it-should-be-derived@cref}{{[originallemma][0][]0DIC}{[1][0][]source}{}{}{}}
\makeatother
\crefname{originallemma}{original lemma}{original lemmas}
\Crefname{originallemma}{Original lemma}{Original lemmas}
\makeatletter
\newlabel{perfect-lemma-flat-proper-pseudo-coherent-direct-image-general}{{0CSD}{source}{}{sourcepointer-3}{}}
\newlabel{perfect-lemma-flat-proper-pseudo-coherent-direct-image-general@cref}{{[originallemma][0][]0CSD}{[1][0][]source}{}{}{}}
\makeatother
\crefname{originallemma}{original lemma}{original lemmas}
\Crefname{originallemma}{Original lemma}{Original lemmas}
\makeatletter
\newlabel{more-algebra-lemma-Rlim-pseudo-coherent-gives-pseudo-coherent}{{0CQF}{source}{}{sourcepointer-4}{}}
\newlabel{more-algebra-lemma-Rlim-pseudo-coherent-gives-pseudo-coherent@cref}{{[originallemma][0][]0CQF}{[1][0][]source}{}{}{}}
\makeatother
\title{Global sections of the twisted object are pseudo-coherent and compute the reductions in the derived existence situation}
\author{The Stacks Project Authors}
\date{stacks-project:flat.tex:local-snapshot-20261011; GNU FDL 1.2 or later}
\begin{document}
\maketitle
\noindent\textit{Editorial note.} All delivered mathematical text is the author's own, copied from the pinned Stacks Project chapter file flat.tex; no statement or proof step was rewritten, reordered or removed.\par\medskip
\noindent\textit{Editorial note.} The preamble is the chapter's own retained minimal preamble (sources/stacks-flat/retained-preamble.tex), a verbatim excerpt of the native Xy-pic setup and of the native theorem declarations and mathematical macros, followed by the project own definitions of the reference, slogan and history display environments (verbatim package plus three newenvironment lines) that the retained file carries at its end and that the project itself typesets as comments. No macro was invented or redefined, and the wrapper class and the amsmath/amssymb/amsthm/graphicx base packages are supplied by the assembler.\par\medskip
\noindent\textit{Editorial note.} The counted statement and proof contain no cross reference, no citation and no included file, so the delivered body is self contained.\par\medskip
\noindent\textit{Editorial note.} Canonical location of the counted proposition: Stacks tag 0DID (sources/stacks-flat/pinned-original-source.tex line 8462).\par\medskip
\noindent\textit{Allowed standard prerequisites (editorial).}\begin{itemize}
\item The object K in the derived existence situation. Same chapter, Lemma labelled lemma-compute-what-it-should-be-derived, invoked in the statement, and Situation labelled situation-existence-derived. Supplies K and the system K\_n.
\item Pseudo-coherence and base change of the finite level global sections. Derived Categories of Schemes, Lemma labelled perfect-lemma-flat-proper-pseudo-coherent-direct-image-general, cited at proof lines 8478-8482. Shows each M\_n is pseudo-coherent over D(A\_n) with compatible base change.
\item Derived limits of pseudo-coherent complexes. More on Algebra, Lemma labelled more-algebra-lemma-Rlim-pseudo-coherent-gives-pseudo-coherent, cited at proof lines 8484-8488. Makes Rlim M\_n pseudo-coherent and recovers its base changes.
\item Commutation of global sections with derived limits. The exactness of RGamma(X, -) on D\_QCoh and its commutation with products, used at proof lines 8488-8500. Identifies RGamma(X, E otimes\textasciicircum{}L K) with Rlim M\_n.
\end{itemize}\medskip
\ifdefined\hypertarget\hypertarget{sourcepointer-1}{}\fi
\noindent\textit{Source reference.} Original reference 0DIB: sources/stacks-flat/pinned-original-source.tex, source lines 8379--8379. Same-chapter pointer into the author environment situation at source line 8379; invoked at proof line None of the counted proof. It is retained only as a navigation pointer; every load-bearing step is inside the delivered span.\par\medskip
\ifdefined\hypertarget\hypertarget{sourcepointer-2}{}\fi
\noindent\textit{Source reference.} Original lemma 0DIC: sources/stacks-flat/pinned-original-source.tex, source lines 8406--8406. Same-chapter pointer into the author environment lemma at source line 8406; invoked at proof line None of the counted proof. It is retained only as a navigation pointer; every load-bearing step is inside the delivered span.\par\medskip
\ifdefined\hypertarget\hypertarget{sourcepointer-3}{}\fi
\noindent\textit{Source reference.} Original lemma 0CSD: sources/stacks-chow/tags.txt, source lines 16028--16028. Cross-chapter standard prerequisite located in the project tag list staged in the local snapshot; invoked at proof line 8480 of the counted proof. It is retained only as a navigation pointer; every load-bearing step is inside the delivered span.\par\medskip
\ifdefined\hypertarget\hypertarget{sourcepointer-4}{}\fi
\noindent\textit{Source reference.} Original lemma 0CQF: sources/stacks-chow/tags.txt, source lines 15960--15960. Cross-chapter standard prerequisite located in the project tag list staged in the local snapshot; invoked at proof line 8486 of the counted proof. It is retained only as a navigation pointer; every load-bearing step is inside the delivered span.\par\medskip
\begin{lemma}
\label{lemma-compute-against-perfect-derived}
In Situation \ref{situation-existence-derived} let $K$ be as in
Lemma \ref{lemma-compute-what-it-should-be-derived}. For any perfect
object $E$ of $D(\mathcal{O}_X)$ the cohomology
$$
M = R\Gamma(X, K \otimes^\mathbf{L} E)
$$
is a pseudo-coherent object of $D(A)$ and there is a canonical isomorphism
$$
R\Gamma(X_n, K_n \otimes^\mathbf{L} E|_{X_n}) = M \otimes_A^\mathbf{L} A_n
$$
in $D(A_n)$. Here $E|_{X_n}$ denotes the derived pullback of $E$ to $X_n$.
\end{lemma}

\begin{proof}
Write $E_n = E|_{X_n}$ and
$M_n = R\Gamma(X_n, K_n \otimes^\mathbf{L} E|_{X_n})$.
By Derived Categories of Schemes, Lemma
\ref{perfect-lemma-flat-proper-pseudo-coherent-direct-image-general}
we see that $M_n$ is a pseudo-coherent object of $D(A_n)$
whose formation commutes with base change. Thus the maps
$M_n \otimes_{A_n}^\mathbf{L} A_{n - 1} \to M_{n - 1}$
coming from $\varphi_n$ are isomorphisms. By
More on Algebra, Lemma
\ref{more-algebra-lemma-Rlim-pseudo-coherent-gives-pseudo-coherent}
we find that $R\lim M_n$ is pseudo-coherent and
that its base change back to $A_n$ recovers $M_n$.
On the other hand, the exact functor
$R\Gamma(X, -) : D_\QCoh(\mathcal{O}_X) \to D(A)$
of triangulated categories commutes with products
and hence with derived limits, whence
$$
R\Gamma(X, E \otimes^\mathbf{L} K) =
R\lim R\Gamma(X,  E \otimes^\mathbf{L} K_n) =
R\lim R\Gamma(X_n, E_n \otimes^\mathbf{L} K_n) =
R\lim M_n
$$
as desired.
\end{proof}

\end{document}
